% from ega2-1-fr.tex
\oldpage[II]{16}
En vertu de (\hyperref[II.1.7.2-fr]{1.7.2}) et de la permutabilité de
$\mathbf{S}$ avec les limites inductives, on a
$(\mathbf{S}(\mathcal{E}))_s=\mathbf{S}(\mathcal{E}_s)$ pour tout point
$s\in S$. Si $\mathcal{E}$, $\mathcal{F}$ sont deux $\mathcal{A}$-Modules,
$\mathbf{S}(\mathcal{E}\oplus\mathcal{F})$ s'identifie canoniquement à
$\mathbf{S}(\mathcal{E})\otimes_{\mathcal{A}}\mathbf{S}(\mathcal{F})$,
comme on le voit sur les préfaisceaux correspondants.

Notons aussi que $\mathbf{S}(\mathcal{E})$ est une $\mathcal{A}$-Algèbre
graduée, somme directe infinie des $\mathbf{S}_n(\mathcal{E})$, où le
$\mathcal{A}$-Module $\mathbf{S}_n(\mathcal{E})$ est le faisceau associé au
préfaisceau $U\mapsto\mathbf{S}_n(\Gamma(U,\mathcal{E}))$. Si on prend en
particulier $\mathcal{E}=\mathcal{A}$, on voit que
$\mathbf{S}_{\mathcal{A}}(\mathcal{A})$ s'identifie à
$\mathcal{A}[T]=\mathcal{A}\otimes_{\mathbf{Z}}\mathbf{Z}[T]$
($T$ indéterminée, $\mathbf{Z}$ étant considéré comme faisceau simple).
\end{env}

\begin{env}[1.7.5]
\label{II.1.7.5-fr}
Soient $(T,\mathcal{B})$ un second espace annelé, $f$ un morphisme
$(S,\mathcal{A})\to(T,\mathcal{B})$. Si $\mathcal{F}$ est un
$\mathcal{B}$-Module, $\mathbf{S}(f^*(\mathcal{F}))$ s'identifie
canoniquement à $f^*(\mathbf{S}(\mathcal{F}))$~; en effet, si
$f=(\psi,\theta)$, on a par définition
(\hyperref[0.4.3.1-fr]{0, 4.3.1})
\[
  \mathbf{S}(f^*(\mathcal{F}))
  =\mathbf{S}(\psi^*(\mathcal{F})
    \otimes_{\psi^*(\mathcal{B})}\mathcal{A})
  =\mathbf{S}(\psi^*(\mathcal{F}))
    \otimes_{\psi^*(\mathcal{B})}\mathcal{A}
\]
(\hyperref[II.1.7.2-fr]{1.7.2})~; pour tout ouvert $U$ de $S$ et toute
section $h$ de $\mathbf{S}(\psi^*(\mathcal{F}))$ au-dessus de $U$, $h$
coïncide, dans un voisinage $V$ de chaque point $s\in U$, avec un élément de
$\mathbf{S}(\Gamma(V,\psi^*(\mathcal{F})))$~; si on se reporte à la
définition de $\psi^*(\mathcal{F})$
(\hyperref[0.3.7.1-fr]{0, 3.7.1}) et qu'on tient compte de ce que tout élément
de $\mathbf{S}(E)$ pour un module $E$ est combinaison linéaire d'un nombre
fini de produits d'éléments de $E$, on voit qu'il y a un voisinage $W$ de
$\psi(s)$ dans $T$, une section $h'$ de $\mathbf{S}(\mathcal{F})$ au-dessus
de $W$ et un voisinage $V'\subset V\cap\psi^{-1}(W)$ de $s$ tels que $h$
coïncide avec $t\mapsto h'(\psi(t))$ au-dessus de $V'$~; d'où notre
assertion.
\end{env}

\begin{proposition}[1.7.6]
\label{II.1.7.6-fr}
Soient $A$ un anneau, $S=\operatorname{Spec}(A)$ son spectre premier,
$\mathcal{E}=\widetilde{M}$ le $\mathcal{O}_S$-Module associé à un $A$-module
$M$~; alors la $\mathcal{O}_S$-Algèbre $\mathbf{S}(\mathcal{E})$ est associée
à la $A$-algèbre $\mathbf{S}(M)$.
\end{proposition}

\begin{proof}
En effet, pour tout $f\in A$,
$\mathbf{S}(M_f)=(\mathbf{S}(M))_f$
(\hyperref[II.1.7.3-fr]{1.7.3}) et la proposition résulte donc des définitions
(\hyperref[I.1.3.4-fr]{I, 1.3.4}).
\end{proof}

\begin{corollary}[1.7.7]
\label{II.1.7.7-fr}
Si $S$ est un préschéma, $\mathcal{E}$ un $\mathcal{O}_S$-Module
quasi-cohérent, la $\mathcal{O}_S$-Algèbre $\mathbf{S}(\mathcal{E})$ est
quasi-cohérente. Si en outre $\mathcal{E}$ est de type fini, chacun des
$\mathcal{O}_S$-Modules $\mathbf{S}_n(\mathcal{E})$ est de type fini.
\end{corollary}

\begin{proof}
La première assertion est conséquence immédiate de
(\hyperref[II.1.7.6-fr]{1.7.6}) et de
(\hyperref[I.1.4.1-fr]{I, 1.4.1})~; la seconde résulte de ce que, si $E$ est
un $A$-module de type fini, $\mathbf{S}_n(E)$ est un $A$-module de type
fini~; on applique alors (\hyperref[I.1.3.13-fr]{I, 1.3.13}).
\end{proof}

\begin{definition}[1.7.8]
\label{II.1.7.8-fr}
Soit $\mathcal{E}$ un $\mathcal{O}_S$-Module quasi-cohérent. On appelle
\emph{fibré vectoriel sur $S$ défini par $\mathcal{E}$} et on note
$\mathbf{V}(\mathcal{E})$ le spectre
(\hyperref[II.1.3.1-fr]{1.3.1}) de la $\mathcal{O}_S$-Algèbre quasi-cohérente
$\mathbf{S}(\mathcal{E})$.
\end{definition}

En vertu de (\hyperref[II.1.2.7-fr]{1.2.7}), pour tout $S$-préschéma $X$, il
y a correspondance biunivoque canonique entre les $S$-morphismes
$X\to\mathbf{V}(\mathcal{E})$ et les homomorphismes de
$\mathcal{O}_S$-Algèbres $\mathbf{S}(\mathcal{E})\to\mathcal{A}(X)$, donc
aussi entre ces $S$-morphismes et les homomorphismes de
$\mathcal{O}_S$-Modules
$\mathcal{E}\to\mathcal{A}(X)=f_*(\mathcal{O}_X)$ (où $f$ est le morphisme
structural $X\to S$). En particulier~:

\begin{env}[1.7.9]
\label{II.1.7.9-fr}
Prenons pour $X$ un sous-préschéma induit par $S$ sur un ouvert $U\subset S$.
Alors, les $S$-morphismes $U\to\mathbf{V}(\mathcal{E})$ ne sont autres que
les $U$-sections (\hyperref[I.2.5.5-fr]{I, 2.5.5}) du $U$-préschéma induit
par $\mathbf{V}(\mathcal{E})$ sur l'ouvert $p^{-1}(U)$ (où $p$ est le
morphisme structural $\mathbf{V}(\mathcal{E})\to S$). D'après ce qu'on vient
de voir, ces $U$-sections correspondent biunivoquement aux homomorphismes de
$\mathcal{O}_S$-Modules
$\mathcal{E}\to j_*(\mathcal{O}_S|U)$ (où $j$ est l'injection canonique
$U\to S$), ou,

